\documentclass[10pt]{article}

\usepackage[margin=0.72in]{geometry}
\usepackage{amsmath,amssymb}
\usepackage{graphicx}
\usepackage{pdfpages}
\usepackage{booktabs}
\usepackage{tabularx}
\usepackage{array}
\usepackage{longtable}
\usepackage{hyperref}
\usepackage{xcolor}
\usepackage{float}
\usepackage{caption}
\usepackage{enumitem}

\hypersetup{
  colorlinks=true,
  linkcolor=blue!55!black,
  urlcolor=blue!55!black
}

\setlength{\parindent}{0pt}
\setlength{\parskip}{0.45em}
\setlist[itemize]{leftmargin=1.3em, itemsep=0.15em, topsep=0.15em}
\captionsetup{font=small, labelfont=bf}

\newcommand{\code}[1]{\texttt{#1}}
\newcommand{\pct}{\%}
\newcommand{\tightsection}[1]{\section{#1}\vspace{-0.35em}}
\newcommand{\nzvec}[6]{[#1, #2, #3, #4, #5, #6]}

\title{Orthogonal Sparsity Pressure: Status Update}
\author{Thiago Mattar\\\href{mailto:thiagocmattar@gmail.com}{thiagocmattar@gmail.com}}
\date{2026-07-05, updated 2026-07-06}

\begin{document}
\maketitle
\vspace{-1.4em}

\tightsection{Immediate Question}

This update counts only completed one-pass MiniPile runs with \code{metrics.json}, \code{manifest.json}, \code{predictions.jsonl}, and a final checkpoint. Pythia-14M is used as a random-initialized architecture, not released checkpoint weights.

\begin{quote}
If MLP activation pressure changes MLP activation and weight distributions, is the pressure also reshaping earlier transformer computations? If so, can we control the site where sparsity appears, and is the poor GELU sparsity/performance tradeoff partly an activation-function problem?
\end{quote}

\tightsection{Hypothesis: MLP Pressure Couples Back Into Earlier Computation}

The first clue came from the MLP-only high-pressure GELU runs. Pressure was trained only on \code{mlp\_hiddens} for RN, OR, L1N, and OL1, but final-checkpoint validation diagnostics show that residual streams and attention outputs also move substantially relative to AdamW. This is evidence of coupling or leakage through the network, although it is not evidence that off-target sites become sparse in the same direction as the pressured MLP site.

\begin{longtable}{p{0.11\textwidth}p{0.25\textwidth}p{0.25\textwidth}p{0.25\textwidth}}
\caption{Validation-cache near-zero mass table for GELU MLP-only high-pressure runs. Entries are \(|a|\le0.01\) percentages by layer [0--5].}\\
\toprule
Method & MLP hiddens & Residual streams & Attention outputs \\
\midrule
\endfirsthead
\toprule
Method & MLP hiddens & Residual streams & Attention outputs \\
\midrule
\endhead
AdamW & \nzvec{7.1}{6.7}{6.7}{6.0}{5.2}{4.7} & \nzvec{36.6}{13.6}{10.7}{11.3}{11.2}{10.4} & \nzvec{61.7}{31.1}{17.0}{25.2}{44.2}{30.2} \\
RN & \nzvec{62.2}{82.6}{77.6}{75.5}{76.6}{72.3} & \nzvec{38.3}{0.0}{0.0}{0.6}{1.2}{0.0} & \nzvec{0.0}{2.9}{6.2}{8.9}{3.0}{0.0} \\
OR & \nzvec{28.5}{75.6}{80.0}{75.2}{71.2}{66.3} & \nzvec{37.8}{2.9}{0.1}{1.8}{0.9}{0.0} & \nzvec{4.1}{3.6}{7.4}{10.6}{14.8}{0.7} \\
L1N & \nzvec{39.0}{74.3}{96.8}{98.7}{98.0}{92.6} & \nzvec{39.6}{0.0}{0.0}{0.0}{0.0}{0.0} & \nzvec{0.0}{0.0}{2.8}{11.4}{2.1}{0.0} \\
OL1 & \nzvec{27.3}{59.6}{88.6}{96.6}{96.1}{91.6} & \nzvec{38.4}{4.1}{0.0}{0.0}{0.0}{0.0} & \nzvec{3.6}{0.2}{1.6}{11.3}{8.7}{0.0} \\
\bottomrule
\end{longtable}

\textbf{Readout.} MLP-only pressure creates very high MLP near-zero mass, especially for L1N/OL1. But off-target residual and attention near-zero mass often decreases relative to AdamW, except for the first residual layer. So the simple story is not ``MLP pressure sparsifies everything''. The more accurate story is: MLP pressure strongly reshapes upstream transformer distributions, and those shifts are visible in residual-stream and attention-output diagnostics.

\begin{figure}[H]
\centering
\includegraphics[width=\textwidth,height=0.60\textheight,keepaspectratio]{../../figures/66-status-update-gelu-mlp-only-coupling-density-comparison.pdf}
\caption{Compact activation-density coupling diagnostic for GELU MLP-only pressure, with log-scaled probability density. Rows compare AdamW, L1N \(w=5\), and RN \(w=1,c=0.05,\sigma=0.05\); columns show MLP hiddens, residual streams, and attention outputs at the same representative layer.}
\end{figure}

\begin{figure}[H]
\centering
\includegraphics[width=\textwidth,height=0.58\textheight,keepaspectratio]{../../figures/70-status-update-gelu-mlp-only-pressure-weight-diagnostic.pdf}
\caption{Compact pressure-weight diagnostic for the same GELU MLP-only comparison. Rows compare AdamW, L1N, and RN; columns aggregate final-checkpoint weights across all six layers for MLP-hidden input weights, residual-stream output-projection weights, and attention QKV weights.}
\end{figure}

\tightsection{Applying Pressure to More Sites}

The coupling hypothesis led to two site-scope follow-ups:
\begin{itemize}
  \item \textbf{All-site pressure}: \code{mlp\_hiddens}, \code{residual\_streams}, and \code{attention\_outputs}.
  \item \textbf{MLP+residual pressure}: \code{mlp\_hiddens} and \code{residual\_streams}, leaving attention outputs unpressured.
\end{itemize}

\subsection*{All-site pressure}

All-site pressure answers one question clearly: attention outputs can be made extremely near-zero under this pressure family. But that happens while MLP hiddens become less sparse than in the MLP-only high-pressure runs. In other words, attention sparsification appears to compete with, or at least redistribute pressure away from, MLP sparsification.

\begin{longtable}{p{0.11\textwidth}p{0.25\textwidth}p{0.25\textwidth}p{0.25\textwidth}}
\caption{Validation-cache near-zero mass table for GELU all-site pressure. Entries are \(|a|\le0.01\) percentages by layer [0--5].}\\
\toprule
Method & MLP hiddens & Residual streams & Attention outputs \\
\midrule
\endfirsthead
\toprule
Method & MLP hiddens & Residual streams & Attention outputs \\
\midrule
\endhead
AdamW & \nzvec{7.1}{6.7}{6.7}{6.0}{5.2}{4.7} & \nzvec{36.6}{13.6}{10.7}{11.3}{11.2}{10.4} & \nzvec{61.7}{31.1}{17.0}{25.2}{44.2}{30.2} \\
OR all-site & \nzvec{21.9}{19.6}{37.9}{47.6}{49.8}{45.7} & \nzvec{43.5}{46.8}{55.7}{59.8}{37.1}{14.4} & \nzvec{98.9}{97.8}{94.2}{94.9}{81.3}{77.8} \\
OL1 all-site & \nzvec{17.6}{16.9}{23.6}{23.8}{18.7}{13.7} & \nzvec{40.9}{29.4}{35.0}{37.1}{42.6}{41.0} & \nzvec{99.3}{98.9}{91.3}{83.4}{93.7}{69.9} \\
\bottomrule
\end{longtable}

\begin{table}[H]
\centering
\caption{All-site endpoint comparison against matched MLP-only high-pressure runs. Validation loss is exact from training metrics; site aggregate near-zero values are the validation-cache diagnostics used in the table above.}
\small
\begin{tabularx}{\textwidth}{llllrrrX}
\toprule
Config & Sites & Method & Setting & Val. loss & MLP agg. & Attn agg. & Readout \\
\midrule
56 & MLP & OR & \(w=1,c=0.05,\sigma=0.05\) & 5.2281 & 66.1\pct & 6.9\pct & MLP sparsity high, validation loss poor. \\
65 & All & OR & \(w=1,c=0.05,\sigma=0.05\) & 5.1033 & 37.1\pct & 90.8\pct & Attention becomes sparse; MLP sparsity drops. \\
59 & MLP & OL1 & \(w=5\) & 5.2342 & 76.6\pct & 4.2\pct & MLP sparsity high, validation loss poor. \\
66 & All & OL1 & \(w=5\) & 4.9192 & 19.0\pct & 89.4\pct & Much better loss, but sparsity moves mainly to attention outputs and residual streams. \\
\bottomrule
\end{tabularx}
\end{table}

\begin{figure}[H]
\centering
\includegraphics[width=\textwidth]{../../figures/67-status-update-all-site-pressure-site-clipping-frontiers.pdf}
\caption{All-site pressure post-hoc clipping frontiers in one comparable view. Each panel clips only one activation family; attention outputs can be pushed to high exact-zero sparsity, while MLP-hidden sparsity remains the expensive axis.}
\end{figure}

\begin{figure}[p]
\centering
\includegraphics[width=\textwidth,height=0.88\textheight,keepaspectratio]{../../figures/41-pythia-14m-minipile-full-pass-all-site-pressure-learning-curves.pdf}
\caption{All-site pressure learning curves. OL1 all-site narrows the validation-loss gap relative to high-pressure MLP-only OL1.}
\end{figure}
\clearpage

\subsection*{MLP+residual pressure}

The MLP+residual runs test whether adding residual streams, but not attention outputs, improves the tradeoff. It does improve the validation loss relative to high-pressure MLP-only L1/OL1, but it again changes where near-zero mass appears. L1N MLP+residual is close to AdamW in validation loss, yet its MLP near-zero mass is much lower than high-pressure MLP-only L1N. Attention outputs also move even without being directly pressured.

\begin{longtable}{p{0.11\textwidth}p{0.25\textwidth}p{0.25\textwidth}p{0.25\textwidth}}
\caption{Validation-cache near-zero mass table for GELU MLP+residual pressure. Entries are \(|a|\le0.01\) percentages by layer [0--5].}\\
\toprule
Method & MLP hiddens & Residual streams & Attention outputs \\
\midrule
\endfirsthead
\toprule
Method & MLP hiddens & Residual streams & Attention outputs \\
\midrule
\endhead
AdamW & \nzvec{7.1}{6.7}{6.7}{6.0}{5.2}{4.7} & \nzvec{36.6}{13.6}{10.7}{11.3}{11.2}{10.4} & \nzvec{61.7}{31.1}{17.0}{25.2}{44.2}{30.2} \\
RN & \nzvec{22.6}{58.4}{64.9}{69.7}{67.5}{55.7} & \nzvec{55.5}{62.9}{61.3}{54.4}{43.7}{22.3} & \nzvec{69.3}{39.1}{15.2}{21.1}{12.9}{5.2} \\
OR & \nzvec{24.0}{34.7}{67.0}{69.9}{63.6}{58.2} & \nzvec{45.8}{53.9}{59.9}{49.9}{39.1}{18.5} & \nzvec{91.1}{50.5}{25.2}{9.1}{13.6}{6.7} \\
L1N & \nzvec{12.7}{12.8}{16.3}{13.2}{14.3}{13.3} & \nzvec{43.5}{37.2}{37.2}{37.4}{39.6}{37.5} & \nzvec{89.2}{45.9}{37.4}{56.9}{63.3}{34.6} \\
OL1 & \nzvec{19.9}{24.2}{33.5}{37.9}{41.5}{20.4} & \nzvec{41.9}{33.2}{43.2}{39.2}{58.2}{68.5} & \nzvec{66.6}{37.7}{36.7}{49.1}{63.6}{46.7} \\
\bottomrule
\end{longtable}

\begin{table}[H]
\centering
\caption{MLP+residual endpoint comparison against MLP-only high-pressure runs. Validation loss is exact from training metrics; site aggregate near-zero values are the validation-cache diagnostics used in the table above.}
\small
\begin{tabularx}{\textwidth}{llllrrrX}
\toprule
Config & Sites & Method & Setting & Val. loss & MLP agg. & Resid. agg. & Readout \\
\midrule
57 & MLP & RN & \(w=1,c=0.05,\sigma=0.05\) & 5.2403 & 74.4\pct & 6.7\pct & High MLP sparsity, poor loss. \\
70 & MLP+resid & RN & \(w=1,c=0.05,\sigma=0.05\) & 5.1969 & 56.5\pct & 50.0\pct & Better loss, residual near-zero mass rises. \\
56 & MLP & OR & \(w=1,c=0.05,\sigma=0.05\) & 5.2281 & 66.1\pct & 7.2\pct & High MLP sparsity, poor loss. \\
71 & MLP+resid & OR & \(w=1,c=0.05,\sigma=0.05\) & 5.1788 & 52.9\pct & 44.5\pct & Better loss, residual site absorbs pressure. \\
58 & MLP & L1N & \(w=5\) & 5.2144 & 83.2\pct & 6.6\pct & Very sparse MLP, poor loss. \\
72 & MLP+resid & L1N & \(w=5\) & 4.8848 & 13.8\pct & 38.7\pct & Near-AdamW loss, but MLP sparsity largely blocked. \\
59 & MLP & OL1 & \(w=5\) & 5.2342 & 76.6\pct & 7.1\pct & Very sparse MLP, poor loss. \\
73 & MLP+resid & OL1 & \(w=5\) & 5.0345 & 29.6\pct & 47.4\pct & Better than MLP-only OL1, still not competitive with AdamW. \\
\bottomrule
\end{tabularx}
\end{table}

\begin{figure}[p]
\centering
\includegraphics[width=\textwidth,height=0.88\textheight,keepaspectratio]{../../figures/52-pythia-14m-minipile-full-pass-mlp-residual-pressure-learning-curves.pdf}
\caption{MLP+residual pressure learning curves for completed configs 70--73. L1N MLP+residual is close to AdamW, but it no longer creates the high MLP near-zero regime seen in high-pressure MLP-only L1N.}
\end{figure}
\clearpage

\begin{figure}[H]
\centering
\includegraphics[width=\textwidth]{../../figures/68-status-update-mlp-residual-pressure-site-clipping-frontiers.pdf}
\caption{MLP+residual pressure post-hoc clipping frontiers in one comparable view, including the AdamW baseline. The site-wise view makes the tradeoff clearer: the best-loss L1N run no longer occupies the high MLP-sparsity regime, while residual and attention robustness differ by method.}
\end{figure}

\tightsection{Is This a GELU Problem? Early ReLU Result}

The site-scope experiments raised a sharper hypothesis: perhaps the GELU MLP nonlinearity itself makes high MLP activation sparsity expensive. We therefore changed the Pythia MLP activation to ReLU while keeping the one-pass MiniPile recipe fixed. The completed ReLU runs so far are AdamW, RN \(w=1,c=0.05,\sigma=0.05\), and OR \(w=1,c=0.05,\sigma=0.05\). ReLU L1N and OL1 are still pending and must not be counted yet.

\begin{table}[H]
\centering
\caption{GELU versus ReLU endpoint comparison for completed AdamW/RN/OR full-pass runs. Validation loss and MLP exact-zero/near-zero endpoint summaries are exact values from each run's training \code{metrics.json}.}
\small
\begin{tabularx}{\textwidth}{llrrrrX}
\toprule
Activation & Method & Val. loss & \(\Delta\) vs same-act AdamW & Exact zero & NZ \(\le0.01\) & Readout \\
\midrule
GELU & AdamW & 4.8317 & 0.0000 & 0.00\pct & 6.03\pct & Baseline full-pass GELU run. \\
GELU & RN & 5.2403 & +0.4085 & 0.00\pct & 72.73\pct & Sparse-ish near zero, but high loss. \\
GELU & OR & 5.2281 & +0.3964 & 0.00\pct & 63.81\pct & Similar poor loss tradeoff. \\
ReLU & AdamW & 4.8404 & 0.0000 & 52.35\pct & 53.88\pct & ReLU baseline already has exact zeros. \\
ReLU & RN & 4.8134 & -0.0271 & 93.91\pct & 94.20\pct & Very sparse MLP hiddens with AdamW-level validation loss. \\
ReLU & OR & 4.8087 & -0.0317 & 93.23\pct & 93.58\pct & Same qualitative result as RN. \\
\bottomrule
\end{tabularx}
\end{table}

\begin{longtable}{p{0.11\textwidth}p{0.25\textwidth}p{0.25\textwidth}p{0.25\textwidth}}
\caption{Validation-cache near-zero mass table for completed ReLU runs. Entries are \(|a|\le0.01\) percentages by layer [0--5].}\\
\toprule
Method & MLP hiddens & Residual streams & Attention outputs \\
\midrule
\endfirsthead
\toprule
Method & MLP hiddens & Residual streams & Attention outputs \\
\midrule
\endhead
AdamW & \nzvec{55.2}{58.9}{52.0}{51.7}{52.6}{51.2} & \nzvec{37.0}{13.0}{9.2}{8.9}{8.8}{8.0} & \nzvec{93.0}{37.4}{18.4}{30.2}{26.3}{19.0} \\
RN & \nzvec{78.1}{97.3}{96.7}{97.6}{97.5}{96.9} & \nzvec{37.6}{14.1}{13.8}{11.9}{12.7}{11.6} & \nzvec{43.6}{21.4}{28.6}{24.1}{25.9}{18.9} \\
OR & \nzvec{74.7}{96.5}{97.8}{97.7}{97.7}{96.9} & \nzvec{37.5}{15.8}{12.6}{13.5}{13.3}{11.5} & \nzvec{78.6}{27.8}{20.5}{33.2}{27.1}{18.9} \\
\bottomrule
\end{longtable}

\textbf{Current interpretation.} This is the most interesting result in the update. For the completed RN/OR settings, ReLU moves the model to roughly 93--94\pct{} exact-zero MLP hidden activation mass with no observed validation-loss penalty, and with point estimates slightly below the ReLU AdamW baseline. That is qualitatively different from GELU RN/OR, where similar pressure settings paid about 0.40 validation-loss points.

\textbf{TODO / pending experiments.} Config 80 (ReLU L1N \(w=5\)) has partial failed attempts and must be rerun from scratch. Config 81 (ReLU OL1 \(w=5\)) has not completed. The ReLU story should be treated as an early result until those runs finish and the key ReLU candidates are repeated across seeds.

\begin{figure}[p]
\centering
\includegraphics[width=\textwidth,height=0.88\textheight,keepaspectratio]{../../figures/59-pythia-14m-minipile-relu-completed-learning-curves.pdf}
\caption{Completed ReLU learning curves for AdamW, RN, and OR. L1N/OL1 are excluded because they did not complete.}
\end{figure}
\clearpage

\begin{figure}[H]
\centering
\includegraphics[width=\textwidth]{../../figures/69-status-update-relu-site-clipping-frontiers.pdf}
\caption{Completed ReLU post-hoc clipping frontiers in one comparable view. The MLP panel is the key activation-function result: RN/OR start near the high exact-zero regime with AdamW-level validation loss, unlike the GELU high-pressure RN/OR tradeoff.}
\end{figure}

\tightsection{Status Takeaways}

\begin{itemize}
  \item MLP-only GELU pressure does not stay local in the representation. Off-target residual and attention distributions change enough that site-specific diagnostics are necessary for every serious comparison.
  \item All-site pressure can drive attention outputs to very high near-zero mass, but this appears to redistribute sparsity away from MLP hiddens.
  \item MLP+residual pressure improves validation loss relative to high-pressure MLP-only L1/Ricker variants, but the best MLP+residual result no longer produces strong MLP sparsity.
  \item ReLU changes the regime: completed ReLU RN/OR runs reach roughly 93--94\pct{} exact-zero MLP hidden activations with AdamW-level validation loss.
  \item The ReLU result is the priority follow-up, but it is incomplete until ReLU L1N/OL1 finish and key runs are repeated across seeds.
\end{itemize}

\end{document}
